%% =================================================================
%%  5.  Results
%% =================================================================
\section{Results}
\label{sec:results}

This section reports what the prototype search engine did and what
it does not establish.  The Lean~4 library proves individual lemmas
that the engine's pruning rules are meant to rely on; for example,
\texttt{rust\_\allowbreak{}sieve\_\allowbreak{}soundness} justifies the
mod-8 rule used by the Rust
\texttt{generate\_\allowbreak{}illegal\_\allowbreak{}z\_\allowbreak{}valuations}
sieve.  No theorem states that the search as a whole is complete, and
the search does not prove any lower bound on quasiperfect numbers.  The
strongest bound we are aware of is $N > 10^{45}$, due to
Alekseyev~\cite{alekseyev2026}; earlier, Hagis and
Cohen~\cite{hagis1982} showed $N > 10^{35}$.


%% -------------------------------------------------------------------
\subsection{Search Coverage and Known Limitations}
\label{subsec:bounds_achieved}

The engine was configured to search up to $10^{\TelemetryMaxLog}$, but
this target is not a result.  The coverage of the search is incomplete
for the following reasons.
\begin{enumerate}
  \item \textbf{Bounded component list.}  Phase~1 only generates prime
    powers $p^{2e}$ with $p < 250{,}000$ and $2e \le 8$
    (\cref{subsec:global_sieve}).  Any QPN with a prime factor of at
    least $250{,}000$, or with a prime appearing to an exponent above~8,
    lies outside the search.
  \item \textbf{Misapplied forced-inclusion rule.}  The cyclotomic
    dependency graph (CDG) ``forced cascade'' treats certain primes as
    components that must appear in~$N$, and prunes a branch when those
    components cannot be added or would push the abundancy past the
    target.  The Lean theorems
    it cites (\texttt{forced\_\allowbreak{}inclusion} and
    \texttt{transitive\_\allowbreak{}forced\_\allowbreak{}inclusion} in
    \texttt{CyclotomicGraph.lean}) show that certain primes divide
    $\sigma(N)$, not that they divide~$N$.  Since
    $\gcd(N, \sigma(N)) = 1$ for a QPN, such primes in fact cannot
    divide~$N$, so this rule has no formal justification and can
    discard valid branches.
  \item \textbf{Pruning bounds relative to the list.}  Several pruning
    bounds, including the suffix-starvation bound, compute the maximum
    achievable abundancy only from components in the engine's finite
    list (\cref{subsec:dfs_prefix}).  They are therefore not valid for
    an arbitrary~$N$ whose factors lie outside that list.
  \item \textbf{Ray-casting did not run.}  The exact ray-casting stage
    (\cref{subsec:ray_casting}) was not exercised; every branch
    recorded in the telemetry was terminated by an earlier pruning
    rule.
\end{enumerate}

For branches whose prefix is disjoint from $\{3, 5\}$ the engine
applies the Prasad--Sunitha bound~\cite{prasad_sunitha}
($\omega(N) \ge \TelemetryPrasadSunithaBound$), which is proved in
Lean.  For prefixes containing both 3 and~5, the bound
$\omega(N) \ge \TelemetryHagisBaselineMinPrimeFactors$ of Hagis and
Cohen is used.

The computation was executed locally.  The hardware environment and physical
execution telemetry metrics are detailed in \cref{tab:hardware_specs}.

\begin{table}[htbp]
    \centering
    \makebox[\textwidth][c]{ \begin{tabular}{|l|l|}
        \hline
        \textbf{Parameter} & \textbf{Specification / Value} \\
        \hline
        Computer / Arch & Apple Mac mini (M4, 10 Cores) \\
        Available RAM & 16 GB Unified Memory \\
        Total Wall-Clock Time & \TelemetryTotalTime \\
        Phase 1 (Precomputation/ECM Factorization) Time & \TelemetryPhaseOneTime \\
        Phase 2 (DFS Traversal) Time & $\approx \TelemetryPhaseTwoTime$ seconds \\
        Phase 1 Eliminations & \TelemetryPhaseOnePruned\ components pruned \\
        Phase 2 Checks & \TelemetryPhaseTwoBranches\ branches evaluated \\
        Run Certificate Hash & \TelemetryCertHash \\
        \hline
    \end{tabular} }
    \caption{Hardware specifications and search telemetry.}
    \label{tab:hardware_specs}
\end{table}

As shown in \cref{tab:hardware_specs}, the main cost
is the up-front Phase~1 integer factorization---specifically computing exact
cyclotomic sums and running the Elliptic Curve Method (ECM) factorization on
tens of thousands of massive $\sigma(p^{2e})$ values before the DFS even starts.
The combinatorial DFS tree traversal (Phase~2) is not the bottleneck; at approximately
\TelemetryNodesPerSec\ nodes per second, it efficiently leverages these precomputations
to process the search space.  Most Phase~2 branches were terminated
by starvation pruning, which, as noted above, is only valid relative to
the engine's component list.  The distribution of branch
terminations is outlined in \cref{tab:pruning_stats}.

\begin{table}[htbp]
  \centering
  \makebox[\columnwidth][c]{ \begin{tabular}{|l|r|}
    \hline
    \textbf{Pruning Mechanism}
      & \textbf{\% of Branches Eliminated} \\
    \hline
    Abundance / Starvation Pruning   & \TelemetryAbundancePct\% \\
    Ray-Casting (Exact Valuations)   &  \TelemetryRaycastPct\% \\
    \hline
  \end{tabular} }
  \caption{Distribution of topological branch pruning strategies across
           DFS execution (Total branches pruned: \TelemetryPruned).}
  \label{tab:pruning_stats}
\end{table}

The telemetry records no runtime panics or overflows during the run.
This shows that the prototype executed without crashing; it does not
show that the search covered every candidate below
$10^{\TelemetryMaxLog}$, for the reasons listed in
\cref{subsec:bounds_achieved}.
